\section{模型分析}\label{sec:analysis}
\subsection{完整候选路径的价值}
表\ref{tab:core}汇总\Model{}与容量匹配上下文模型的配对结果，覆盖ETTh1 $(.10,16)$ 和PAMAP2 $(.15,8)$。两数据集的三对训练种子均支持\Model{}具有更低的最终MSE和MAE。最终MSE分别降低\GainETTMSE{}、\GainPAMMSE{}，MAE分别降低\GainETTMAE{}、\GainPAMMAE{}。粗输出已经具有收益；对同一结构精恢复器独立训练后，差距仍然保留，说明显式Patch证据在容量匹配的上下文恢复之外提供了增量价值。

\begin{table*}[t]
 \centering\small\setlength{\tabcolsep}{4pt}
 \begin{tabular}{llrrr}
\toprule
\Lang{数据集}{Dataset} & \Lang{模型}{Model} & \Lang{粗MSE}{\shortstack{Coarse\\MSE}} & \Lang{精MSE}{\shortstack{Refined\\MSE}} & \Lang{精MAE}{\shortstack{Refined\\MAE}} \\
\midrule
ETTh1 & \Model{} & $0.384873\,\pm\,0.006236$ & $0.334986\,\pm\,0.006672$ & $0.372039\,\pm\,0.004057$ \\
ETTh1 & \Lang{上下文对照}{Context-only} & $0.408733\,\pm\,0.001863$ & $0.353320\,\pm\,0.009562$ & $0.385618\,\pm\,0.010399$ \\
PAMAP2 & \Model{} & $0.336419\,\pm\,0.000680$ & $0.324572\,\pm\,0.001591$ & $0.230532\,\pm\,0.000622$ \\
PAMAP2 & \Lang{上下文对照}{Context-only} & $0.351320\,\pm\,0.002814$ & $0.332533\,\pm\,0.001960$ & $0.234430\,\pm\,0.001031$ \\
\bottomrule
\end{tabular}

 \caption{候选路径配对分析。两个数据集均为三个训练种子，每模型平均二十组共享测试遮挡。PAMAP2机制分析队列与表\ref{tab:formal}的正式结果队列分开；粗精结果来自同一选定模型对。}
 \label{tab:core}
\end{table*}

\subsection{来源组织与逐点读取}
表\ref{tab:components}中，在两个分析数据集上分别移除任一来源，都会增加最终平均MSE。ETTh1中跨变量证据的来源删除效应最大，PAMAP2中时间证据最大，说明来源价值随信号特性而变化。

来源内随机选择与统一检索的平均MSE均高于完整结构，表明收益来自来源相关的组织，而不只是提供同等数量的合法片段。Patch共享读取改变了精度取舍：其ETTh1 MSE更低，而逐点读取具有更低的ETTh1 MAE与PAMAP2 MSE。因此同时使用两种误差考察读取粒度。

\begin{table*}[t]
 \centering\small\setlength{\tabcolsep}{3pt}
 \begin{tabular}{lrrrr}
\toprule
\Lang{变体}{Variant} & ETTh1 MSE & ETTh1 MAE & PAMAP2 MSE & PAMAP2 MAE \\
\midrule
\Model{} & $0.334986\,\pm\,0.006672$ & $0.372039\,\pm\,0.004057$ & $0.322915$ & $0.230183$ \\
\Lang{去时间证据}{w/o temporal evidence} & $0.343612\,\pm\,0.006384$ & $0.381684\,\pm\,0.004346$ & $0.333096$ & $0.235096$ \\
\Lang{去跨变量证据}{w/o variable evidence} & $0.349493\,\pm\,0.018407$ & $0.381209\,\pm\,0.010604$ & $0.325663$ & $0.231132$ \\
\Lang{去历史证据}{w/o historical evidence} & $0.339364\,\pm\,0.004246$ & $0.376169\,\pm\,0.002895$ & $0.324277$ & $0.231290$ \\
\Lang{来源内随机}{Source-random} & $0.336177\,\pm\,0.006092$ & $0.377530\,\pm\,0.004732$ & $0.325863$ & $0.231077$ \\
\Lang{统一检索}{Unified retrieval} & $0.339484\,\pm\,0.005147$ & $0.380305\,\pm\,0.002030$ & $0.331974$ & $0.232928$ \\
\Lang{无候选容量匹配}{Context-only (matched)} & $0.353320\,\pm\,0.009562$ & $0.385618\,\pm\,0.010399$ & $0.334383$ & $0.235483$ \\
\Lang{Patch共享读取}{Patch-shared read} & $0.333639\,\pm\,0.004381$ & $0.374103\,\pm\,0.002174$ & $0.323958$ & $0.229965$ \\
\bottomrule
\end{tabular}

 \caption{精恢复后的组件与检索对照。ETTh1为三个训练种子均值$\pm$样本标准差，PAMAP2组件行统一使用seed2023；每个模型均评价二十组共享遮挡。PAMAP2的三种子\Model{}与容量匹配上下文模型对照单列于表\ref{tab:core}。}
 \label{tab:components}
\end{table*}

\subsection{观测条件化精恢复}
在\Model{}配对队列中，精恢复将ETTh1 MSE从0.384873降至0.334986，PAMAP2从0.336419降至0.324572，分别下降12.96\%与3.52\%。该网络将时间修正与由可见观测锚定的参考信息结合。逐坐标诊断分别统计误差减少和增加的贡献，使总体修正能够追溯到具体位置。

表\ref{tab:backbone}进一步将同一精恢复应用于不同骨干。在BeijingAirQuality$(.20,8)$中，T1的MSE从0.225523变为0.129252，\Model{}从0.179570变为0.123941。公共精恢复解释了原输出差距中的较大部分，同时保留\Model{}粗估计的额外优势，从而区分骨干和精恢复的作用。

\begin{table*}[t]
 \centering\small\setlength{\tabcolsep}{4pt}
 \begin{tabular}{lrrrr}
\toprule
\Lang{骨干}{Backbone} & \Lang{粗MSE}{\shortstack{Coarse\\MSE}} & \Lang{精MSE}{\shortstack{Refined\\MSE}} & \Lang{粗MAE}{\shortstack{Coarse\\MAE}} & \Lang{精MAE}{\shortstack{Refined\\MAE}} \\
\midrule
T1 & $0.225523$ & $0.129252$ & $0.264340$ & $0.193203$ \\
\Lang{上下文对照}{Context-only} & $0.188083$ & $0.124760$ & $0.243393$ & $0.189678$ \\
\Model{} & $0.179570$ & $0.123941$ & $0.235067$ & $0.188432$ \\
\bottomrule
\end{tabular}

 \caption{BeijingAirQuality$(.20,8)$上三个冻结骨干接相同结构、相同预算精恢复的结果，seed2023，二十组测试遮挡。粗精误差在相同位置评价。}
 \label{tab:backbone}
\end{table*}

\subsection{不同缺失条件下的诊断}
按掩码定义的缺口长度、最近可见点距离、Patch可见率和窗口边界分组，分组在读取预测前固定。ETTh1$(.10,16)$中，到可见观测距离17--32及33--95的位置，在三个训练种子上均保持正向收益，支持Patch证据在局部观测稀少时发挥作用。对于所有覆盖Patch均无可见观测的坐标，模型仍可利用Patch外部上下文及其他证据来源恢复。完整分层数值与坐标支撑保存在数值证据表中。

\subsection{精度与成本}
表\ref{tab:efficiency}量化候选路径的额外成本。\Model{}与容量匹配上下文模型有效参数接近，候选构建与读取增加推理延迟和allocated显存。PAMAP2中，两者参数分别为575074和574794，完整流程每batch的p50延迟为98.48和55.53ms，使精度收益能够与实际执行代价一起考察。

\begin{table*}[t]
 \centering\small\setlength{\tabcolsep}{5pt}
 \begin{tabular}{llrrr}
\toprule
\Lang{数据集}{Dataset} & \Lang{模型}{Model} & \Lang{参数}{Params} & p50 (ms) & \Lang{显存}{VRAM} (GiB) \\
\midrule
PAMAP2 & \Model{} & 575,074 & 98.48 & 6.121 \\
PAMAP2 & \Lang{上下文对照}{Context-only} & 574,794 & 55.53 & 4.122 \\
ETTh1 & \Model{} & 353,248 & 17.63 & 1.067 \\
ETTh1 & \Lang{上下文对照}{Context-only} & 352,968 & 11.51 & 0.709 \\
\bottomrule
\end{tabular}

 \caption{选定seed2023模型的推理成本；同一独占RTX5090、FP32、batch128，二十次预热后进行一百次同步测量。显存为峰值allocated。}
 \label{tab:efficiency}
\end{table*}

\paragraph{读取敏感性。}
表\ref{tab:sensitivity}在总宽64不变时改变逐点读取头数。参数量不变，用于观察特征划分方式的影响。主模型保持预定八头配置，敏感性实验使用seed2023及验证选定权重。

\begin{table*}[t]
 \centering\small\setlength{\tabcolsep}{5pt}
 \begin{tabular}{crrrr}
\toprule
\Lang{头数}{Heads} & Val MSE & Val MAE & Test MSE & Test MAE \\
\midrule
4 & $0.422429$ & $0.427810$ & $0.326837$ & $0.370128$ \\
8 & $0.446899$ & $0.436511$ & $0.336727$ & $0.373539$ \\
16 & $0.425259$ & $0.428136$ & $0.325889$ & $0.369931$ \\
\bottomrule
\end{tabular}

 \caption{ETTh1$(.10,16)$上的逐点读取头数敏感性，总宽64，seed2023，每模型二十组共享测试遮挡。}
 \label{tab:sensitivity}
\end{table*}
